\chapter{Network Stack}
\label{ch:network}

Veilbox uses a minimal but functional network stack built on BusyBox
networking utilities and the Linux kernel's networking subsystem. This
chapter covers the network configuration, DHCP, routing, and integration
with container networking.

\section{Network Architecture}

The Veilbox network stack has four layers:

\begin{enumerate}[nosep]
  \item \textbf{Physical/virtual device}: VirtIO-NET (virtio\_net driver)
        provides the network interface to QEMU's virtual switch.
  \item \textbf{Link layer}: The kernel's networking core manages the
        Ethernet device (\texttt{eth0}) via the VirtIO-NET driver.
  \item \textbf{Network layer}: IPv4 protocol stack with ARP, ICMP, and
        routing.
  \item \textbf{Application layer}: udhcpc (DHCP client) and user
        applications (wget, nerdctl image pulls, ssh, etc.).
\end{enumerate}

\section{Interface Configuration}

\subsection{Loopback Interface}

The loopback interface provides local communication:

\begin{lstlisting}[language=bash]
ifconfig lo 127.0.0.1 up
\end{lstlisting}

\subsection{Ethernet Interface}

The primary Ethernet interface (\texttt{eth0}) acquires an IP address
via DHCP:

\begin{lstlisting}[language=bash]
ifconfig eth0 0.0.0.0 up
udhcpc -i eth0 -b
\end{lstlisting}

\section{DHCP Client (udhcpc)}

The BusyBox DHCP client is a minimal implementation of the DHCP protocol:

\subsection{Protocol Flow}

\begin{enumerate}[nosep]
  \item \textbf{DHCPDISCOVER}: Client broadcasts a discovery packet to
        255.255.255.255 (UDP port 67).
  \item \textbf{DHCPOFFER}: Server responds with an offered IP address,
        lease time, subnet mask, gateway, and DNS servers.
  \item \textbf{DHCPREQUEST}: Client requests the offered address.
  \item \textbf{DHCPACK}: Server acknowledges the lease.
  \item \textbf{Bound}: Client configures the interface using the
        supplied parameters.
\end{enumerate}

\subsection{Lease Renewal}

At 50\% of lease time, the client attempts to renew. On success, the
renew script is called. On failure, the client tries again at 87.5\%.

\section{Routing}

Default route configuration:

\begin{lstlisting}[language=bash]
route add default gw 10.0.2.2
\end{lstlisting}

In QEMU SLiRP networking, the gateway is always 10.0.2.2. DNS is
at 10.0.2.3. The guest IP is 10.0.2.15/24.

\section{Container Networking}

\subsection{veth Pairs}

When nerdctl creates a container with network access, it creates a veth
pair — two virtual Ethernet interfaces connected by a virtual cable:

\begin{lstlisting}
# Host side: vethXXXXX in root namespace
# Container side: eth0 in container namespace
\end{lstlisting}

\subsection{Network Namespace Isolation}

Each container receives its own network namespace:

\begin{itemize}[nosep]
  \item \textbf{Default}: A veth pair connects the container to the host's
        network. The container shares the host's IP stack (NAT).
  \item \textbf{Host mode} (\texttt{--net=host}): The container uses the
        host's network namespace directly, sharing the same IP address.
  \item \textbf{None} (\texttt{--net=none}): The container has only a
        loopback interface, completely isolated.
\end{itemize}


\section{Network Interface Naming}

Veilbox uses the traditional kernel naming scheme (\texttt{eth0}) rather
than the predictable naming scheme (\texttt{enp0s3}) used by modern
distributions. This is configured via the kernel command line:

\begin{lstlisting}
net.ifnames=0
\end{lstlisting}

\section{TCP/IP Stack Tuning}

The Veilbox kernel includes the following TCP/IP features:

\begin{lstlisting}
CONFIG_TCP_CONG_CUBIC=y    # Default congestion control
CONFIG_TCP_CONG_BBR=y      # BBR congestion control (optional)
CONFIG_IP_MULTICAST=y      # Multicast support
CONFIG_IP_PNP=y            # IP autoconfiguration (DHCP)
CONFIG_SYN_COOKIES=y       # SYN flood protection
\end{lstlisting}

\subsection{Network Performance Tuning}

Parameters in /proc/sys/net/ for tuning:

\begin{lstlisting}
# Increase TCP buffer sizes
echo 2097152 > /proc/sys/net/core/rmem_max
echo 2097152 > /proc/sys/net/core/wmem_max

# Enable TCP window scaling
echo 1 > /proc/sys/net/ipv4/tcp_window_scaling

# Reduce TIME_WAIT sockets
echo 1 > /proc/sys/net/ipv4/tcp_tw_reuse

# Increase backlog queue
echo 1024 > /proc/sys/net/core/somaxconn
\end{lstlisting}

\section{Stateless Address Autoconfiguration}

In addition to DHCP, the kernel supports IPv6 SLAAC (Stateless Address
Autoconfiguration). While Veilbox primarily uses IPv4, the IPv6 stack
is available:

\begin{lstlisting}
CONFIG_IPV6=y
CONFIG_IPV6_ROUTER_PREF=y
CONFIG_IPV6_OPTIMISTIC_DAD=y
\end{lstlisting}

\section{Conntrack and NAT}

Connection tracking is essential for container NAT:

\begin{lstlisting}
# View connection tracking table
cat /proc/net/nf_conntrack

# Maximum tracked connections
echo 65536 > /proc/sys/net/netfilter/nf_conntrack_max

# Connection timeout (seconds)
echo 600 > /proc/sys/net/netfilter/nf_conntrack_tcp_timeout_established
\end{lstlisting}

\section{DNS Resolution}

DNS is configured via /etc/resolv.conf, written by the DHCP script:

\begin{lstlisting}
# /etc/resolv.conf
nameserver 10.0.2.3
\end{lstlisting}

In QEMU SLiRP networking, the DNS server is always at 10.0.2.3, which
forwards queries to the host's DNS resolver.

\section{Advanced DHCP Configuration}

\subsection{Static IP Configuration}

If DHCP is not available, configure a static IP:

\begin{lstlisting}
# Static IP configuration
ifconfig eth0 192.168.1.100 netmask 255.255.255.0 up
route add default gw 192.168.1.1
echo "nameserver 8.8.8.8" > /etc/resolv.conf
echo "nameserver 1.1.1.1" >> /etc/resolv.conf
\end{lstlisting}

\subsection{DHCP with Custom Script}

The udhcpc default script can be customized:

\begin{lstlisting}[language=bash]
#!/bin/sh
# Custom DHCP script /usr/share/udhcpc/default.script
RESOLV_CONF="/etc/resolv.conf"

case "$1" in
    bound|renew)
        ifconfig $interface $ip netmask $subnet
        route add default gw $router 2>/dev/null

        # Write DNS servers
        > $RESOLV_CONF
        for ns in $dns; do
            echo "nameserver $ns" >> $RESOLV_CONF
        done

        # Set hostname if provided
        [ -n "$hostname" ] && hostname "$hostname"

        # Log the event
        echo "DHCP: $interface -> $ip via $router" \
            >> /var/log/messages
        ;;
    deconfig)
        ifconfig $interface 0.0.0.0
        ;;
    leasefail)
        echo "DHCP: lease failed for $interface" \
            >> /var/log/messages
        ;;
esac
\end{lstlisting}

\section{Network Troubleshooting Tools}

\subsection{Ping and Connectivity}

\begin{lstlisting}
# Test basic connectivity
ping -c 3 10.0.2.2      # Gateway (QEMU SLiRP)
ping -c 3 10.0.2.3      # DNS server
ping -c 3 8.8.8.8       # Internet (if NAT works)

# Test DNS resolution
ping -c 3 google.com
\end{lstlisting}

\subsection{Traceroute}

BusyBox traceroute shows the network path:

\begin{lstlisting}
traceroute 8.8.8.8
\end{lstlisting}

\subsection{Tcpdump}

For packet-level analysis, tcpdump can be installed in a container:

\begin{lstlisting}
nerdctl run --rm --net=host alpine sh -c \
    "apk add tcpdump && tcpdump -i eth0"
\end{lstlisting}
